% NOTE (migration 23 Nov 2026): Text moved unchanged from main.tex. Much of it was drafted with AI assistance; rewrite it in your own words before submission (the supervisor runs AI checks).
% NOTE (23 Nov 2026): Research questions, scope, contributions and organisation updated to match the rewritten Methodology. RQ3/RQ4 reframing assumes the supervisor accepts RQ-EVIDENCE.md; change it together with Methodology Section 3.1 otherwise.
\hypertarget{introduction}{%
\chapter{Introduction}\label{introduction}}


\hypertarget{problem-statement}{%
\section{Problem statement}\label{problem-statement}}

Administrators must record scheduling requirements accurately enough to protect access to required teaching while accommodating staff availability and suitable facilities. If the system misunderstands a requirement, a timetable can appear acceptable while failing the intended institutional need. The practical problem is therefore how to support clear requirement entry, useful clarification, and informed approval without adding excessive administrative effort.

The five reviewed timetabling papers provide evidence about scheduling needs and capabilities, but do not evaluate ordinary-language requirement entry with this review workflow. This thesis investigates that task within a limited timetabling scope. Its gap claim is bounded by the reviewed sources; a wider documented literature search remains planned.

\hypertarget{research-gap}{%
\section{Research gap}\label{research-gap}}

Within the five retained timetabling papers, constraint acquisition is a developer or
administrator task and is not evaluated as a research problem (Chapter~\ref{literature-review}). The additional LLM-related sources verified so far address
different questions: LLMs solving scheduling problems described in text
\cite{sharma2026schedbench}, a pilot on appliance scheduling in Italian
\cite{sanguinetti2025preferences}, and a benchmark of LLMs writing constraint
models from problem descriptions~\cite{michailidis2025cpbench},
whose task differs because the output is a general constraint program
with no institutional data to resolve against and no review step. The
gap addressed here is therefore the \emph{evaluation of
natural-language-to-constraint extraction for university course
timetabling}, with human review and independent validation, using a
fixed formal model so that extraction quality can be separated from
solver behaviour. The novelty claim is limited to this evaluated
combination; the use of an LLM is not itself claimed to be novel.

\hypertarget{research-questions}{%
\section{Research questions}\label{research-questions}}

\begin{itemize}
\tightlist
\item
  \textbf{RQ1 (primary).} How accurately can an LLM extract formal
  university scheduling constraints from natural-language requirements
  --- including entity references, parameters and hard/soft
  classification --- when output is restricted to a closed constraint
  vocabulary?
\item
  \textbf{RQ2.} Can the system identify ambiguous, missing,
  contradictory or unsupported requirements and request useful
  clarification instead of guessing?
\item
  \textbf{RQ3.} Does AI-assisted constraint entry with human review
  reduce entry time and the number of corrections compared with manual
  form entry while preserving correct interpretation?
\item
  \textbf{RQ4.} How does a locally hosted base language model compare
  with the same model fine-tuned on reviewed, guide-grounded examples,
  in answer quality and latency, on the same held-out cases?
\end{itemize}

RQ1 and RQ2 are evaluated overall and separately for English, formal Arabic, Egyptian Arabic and mixed Arabic/English, by comparing a rule-based parser with a locally hosted language model on development data with draft labels. The results are therefore development-level evidence, not a confirmatory held-out result. RQ3 could not be studied, because no participants or ethics approval were available in the project period; it is discussed as future work. RQ4 originally compared several commercial models; because no budget for paid calls was approved, it compares a local base model with its fine-tuned versions (Chapter~\ref{methodology-proposed-study}).

A supporting engineering question concerns whether the surrounding
system can keep the model honest: \emph{can an independent validator and
an approval workflow ensure that only reviewed constraints reach the
solver and that no schedule with a supported hard-rule violation is
published?} This is verified by testing
(Section~\ref{scheduling-validation-and-system-behaviour-evaluation}), not by statistical inference.

\hypertarget{objectives}{%
\section{Objectives}\label{objectives}}

\begin{enumerate}
\def\labelenumi{\arabic{enumi}.}
\tightlist
\item
  Identify and document the supported educational and administrative scheduling requirements.
\item
  Enable administrators to enter requirements in ordinary language, inspect their interpretation, and clarify or correct them before approval.
\item
  Generate candidate course timetables that respect approved attendance, availability and facility requirements.
\item
  Give administrators understandable checks and feedback for generated timetables and manual changes before publication.
\item
  Evaluate interpretation correctness separately from timetable quality, and evaluate the local assistant model before and after fine-tuning.
\item
  Deliver a working prototype, reproducible evaluation records and this thesis.
\end{enumerate}

\hypertarget{contribution-proportional-to-the-evidence}{%
\section{Contribution (proportional to the
evidence)}\label{contribution-proportional-to-the-evidence}}

The contributions are: (i) a documented set of supported rules and a
machine-readable output contract for timetabling requirements, with
explicit decisions for clarification, unsupported and conflicting
requests; (ii) a synthetic development corpus of 144 statements in four
language categories with a versioned annotation guide, whose labels are
draft and unreviewed; (iii) a development-level comparison of a
rule-based parser and a locally hosted language model on that corpus,
with family-level uncertainty; (iv) a fine-tuning study of a local
assistant model in four rounds, including a round that made it worse
and the reason; and (v) a working prototype for teaching and examination
timetables in which an approval workflow and an independent validator
keep unreviewed or invalid output from being used. No claim is made
about accuracy on real administrators' statements or about reduced
administrative effort.

\hypertarget{scope}{%
\section{Scope}\label{scope}}

\textbf{Project scope:} JadwAI supports course timetabling for one semester at a time, using a repeating weekly pattern, from a single department up to a whole university built in parts. Lectures, tutorials and laboratories have durations and weekly occurrences; student groups are the unit of student-conflict checking. Natural-language requirements in English, formal Arabic, Egyptian Arabic and mixed Arabic/English are part of the core research scope, alongside an English/Arabic interface with right-to-left presentation. JAI (pronounced ``jay'') is the website's chat assistant; its language model runs on the same server as the application.

The solver chooses instructors, times and rooms. Eligibility is recorded separately for lecture, tutorial and laboratory activities. Each activity retains one instructor across its weekly occurrences; instructor, time and room assignments may be locked. Mandatory weekly instructor workload limits and daily limits for instructors and groups are supported. The unit is elapsed teaching hours: a 90-minute session contributes 1.5 hours, counted once per instructor irrespective of how many groups attend; fixed breaks are excluded. Configured weekly instructor maxima and global hard daily caps are implemented. Preferred daily-limit variants and entity-specific daily overrides are not implemented.

\textbf{Added during the project:} examination timetabling (quizzes, midterms and finals, with halls and TA invigilators), repair suggestions for impossible timetables (proposed to the administrator, never applied automatically), and generation in parts for university-scale data. \textbf{Deferred extensions:} individual enrolments, minimal-change regeneration, calendar feeds and notifications. \textbf{Excluded:} native mobile applications, attendance, payroll, full student-information-system integration, automated registration or sectioning, and institution-wide production deployment. The implementation status of each requirement is listed in Appendix~\ref{app:requirements}.

\hypertarget{organisation-of-the-thesis}{%
\section{Organisation of the thesis}\label{organisation-of-the-thesis}}

Chapter~\ref{chap:background} introduces university timetabling and the constraint-entry problem. Chapter~\ref{literature-review} reviews five timetabling papers through educational needs, implemented capabilities and reported evidence, and examines a study of natural-language preference interpretation. Chapter~\ref{methodology-proposed-study} describes the method followed: the research design, the interpretation contract, the development corpus, the methods compared, the metrics, the fine-tuning of the assistant model, the engineering verification, and the threats to validity. Chapter~\ref{chap:implementation} describes the implemented system and the JAI assistant. Chapter~\ref{evidence} reports the measurements, and Chapter~\ref{chap:conclusion} summarises the findings, limitations and future work. The appendices list the requirements, screenshots of the prototype, and the project timeline with traceability to the requirements.
